\chapter{Dados, telemetria e monitoramento em tempo real}
\label{cap:dados}

\begin{objetivos}
Ao final deste capítulo você deve ser capaz de:
\begin{itemize}
  \item distinguir dado, informação e indicador, e construir o terceiro a partir dos dois primeiros;
  \item calcular e interpretar os indicadores usuais de produção e de consumo;
  \item decidir onde processar cada tarefa: no CLP, na borda ou na nuvem;
  \item reconhecer os sinais simples de degradação usados em manutenção preditiva;
  \item dizer, com clareza, o que o controlador pode e o que ele não pode fazer.
\end{itemize}
\end{objetivos}

\section{Do dado bruto ao indicador}

Os três termos são usados como sinônimos no dia a dia e não são:

\begin{itemize}
  \item \textbf{dado}: o valor lido, com carimbo de tempo. ``Contador de peças = 4.812, às 14h32'';
  \item \textbf{informação}: o dado com contexto. ``4.812 peças produzidas no turno A, contra 5.100 no
        mesmo turno de ontem'';
  \item \textbf{indicador}: a informação transformada em medida de desempenho, comparável entre
        períodos e máquinas. ``Eficiência global do equipamento: 71\,\%''.
\end{itemize}

Toda a cadeia de monitoramento existe para chegar ao terceiro. Uma planta pode ter milhares de tags
e nenhum indicador --- e, nesse caso, tem telemetria sem gestão.

\begin{figure}[H]
  \centering
  \includegraphics[width=0.74\linewidth]{19-fluxo-dados.pdf}
  \caption{Caminho do dado até a decisão, e o retorno da decisão ao processo.}
  \label{fig:fluxo-dados}
\end{figure}

\begin{atencao}
A etapa mais esquecida da Figura~\ref{fig:fluxo-dados} é a última. Indicador que não volta como
decisão --- mudar um \emph{setpoint}, antecipar uma manutenção, reprogramar um lote --- é relatório
de gaveta. Antes de instrumentar qualquer coisa, escreva qual decisão depende daquele número.
\end{atencao}

\index{indicador de desempenho}\index{telemetria}

\section{Os indicadores usuais}

\begin{table}[H]
  \centering
  \caption{Indicadores usuais de produção e consumo, o que exigem do controle e onde costumam falhar.}
  \label{tab:indicadores}
  \begin{tabular}{p{3.0cm} p{4.1cm} p{3.4cm} p{3.8cm}}
    \toprule
    Indicador & Definição conceitual & O que o CLP fornece & Onde costuma falhar \\
    \midrule
    OEE (eficiência global do equipamento)
      & Disponibilidade $\times$ desempenho $\times$ qualidade
      & Estado da máquina (parada, operando, falha), contagem de peças boas e refugadas, tempo
        disponível
      & Classificação de parada mal definida: toda parada vira ``outros'', e o indicador deixa de
        apontar causa \\
    \addlinespace
    Disponibilidade
      & Tempo em operação dividido pelo tempo programado
      & Sinal de máquina em ciclo, tempo em falha
      & Não distinguir parada planejada de não planejada \\
    \addlinespace
    Tempo de ciclo
      & Tempo entre duas peças conforme, ou por etapa do ciclo
      & Marcações de início e fim, com resolução adequada
      & Medir no supervisório um fenômeno que acontece em milissegundos: a média esconde a variação \\
    \addlinespace
    Contagem de produção
      & Peças ou lotes produzidos por turno, produto e linha
      & Contador de ciclo, com tratamento de primeira varredura e de retomada
      & Retenção de contador perdida na queda de energia, ou contagem dupla na partida \\
    \addlinespace
    Consumo energético
      & Energia por tonelada, por peça ou por turno
      & Leitura de medidor por rede ou por pulso; estado de motores e de iluminação
      & Comparar períodos com produtos diferentes e concluir que a eficiência mudou
        (Capítulo~\ref{cap:eficiencia}) \\
    \bottomrule
  \end{tabular}
\end{table}

Uma observação sobre o OEE: ele é útil exatamente porque separa disponibilidade, desempenho e
qualidade. Quando combinado em um único número sem detalhamento, perde a capacidade de apontar
\emph{onde} intervir. O indicador agregado é para a reunião de gestão; o detalhamento é para a
engenharia.

\section{Onde processar: CLP, borda ou nuvem}

Esta é a decisão de arquitetura que define o custo e a utilidade do sistema.

\begin{table}[H]
  \centering
  \caption{Critérios para decidir onde processar cada tarefa.}
  \label{tab:onde-processar}
  \begin{tabular}{p{2.4cm} p{3.3cm} p{3.3cm} p{5.3cm}}
    \toprule
    Critério & No CLP & Na borda & Na nuvem ou na sala de controle \\
    \midrule
    Prazo
      & Prazo garantido, em milissegundos
      & Responde em centenas de milissegundos a segundos
      & Prazo não garantido; depende de banda e de fila \\
    \addlinespace
    Volume
      & Baixo: o que cabe na memória de trabalho
      & Intermediário: filtra e agrega antes de enviar
      & Alto: meses ou anos de histórico de muitas máquinas \\
    \addlinespace
    Criticidade
      & Só o que não pode parar a máquina
      & O que melhora a operação sem risco imediato
      & O que orienta decisão de dias ou meses \\
    \addlinespace
    Manutenção
      & Exige parada e equipe de automação
      & Atualizável em janela curta
      & Atualizável continuamente, por equipe de TI \\
    \bottomrule
  \end{tabular}
\end{table}

Lendo a Tabela~\ref{tab:onde-processar} de cima para baixo chega-se à regra prática:
\textbf{leitura rápida e intertravamento no controlador; transformação e contexto na borda; acúmulo e
análise pesada fora do chão de fábrica}.

Essa divisão também é uma proteção. Cada função que sai do controlador deixa de depender do
determinismo dele --- e, em troca, ganha capacidade de processamento. Mover a decisão certa para o
lugar errado é a origem tanto de máquinas inseguras quanto de indicadores inúteis.

\index{processamento na borda}\index{nuvem!limite de uso}

\section{Detecção de desvio e manutenção preditiva simples}

Nem toda manutenção preditiva exige aprendizado de máquina. Quatro recursos simples, implementáveis
em lógica de controlador ou em ferramenta de histórico, cobrem boa parte dos casos:

\begin{enumerate}
  \item \textbf{limite absoluto}: o valor passou de um patamar. É o alarme clássico;
  \item \textbf{taxa de mudança}: a variável está subindo mais rápido que o normal. Um aumento lento
        e persistente de corrente em um motor indica desgaste ou obstrução, mesmo dentro da faixa
        ``normal'';
  \item \textbf{contagem de ciclos}: o atuador número $n$ vai superar a vida declarada. Simples,
        verificável e diretamente acionável;
  \item \textbf{desvio do padrão}: o tempo de uma etapa do ciclo cresceu 8\% em relação à referência
        do mesmo produto, sem que nenhum limite tenha sido atingido.
\end{enumerate}

Os dois últimos são os mais valiosos, porque antecipam a falha sem depender de um limite fixo --- que
existe justamente para não ser atingido em operação normal.

\begin{atencao}
Detecção de desvio sem contexto gera alarme falso, e alarme falso ensina o operador a ignorar
alarmes. Antes de publicar um indicador de degradação, defina três coisas: a referência, a banda de
tolerância e o que fazer quando o desvio aparecer. Um indicador que ninguém sabe como usar é ruído
com carimbo de tempo.
\end{atencao}

\section{O limite do controlador}

Vale ser explícito, porque é o erro mais comum em projetos de monitoramento:

\begin{itemize}
  \item o CLP faz \textbf{aquisição rápida}, lógica de \textbf{intertravamento}, contagem e
        temporização com precisão conhecida --- e faz isso sem depender de rede;
  \item o CLP \textbf{não} é o lugar de treinar modelo, cruzar meses de histórico, gerar relatório
        gráfico ou calcular indicador que depende de dados de várias máquinas. Havia memória
        contada para isso no Capítulo~\ref{cap:memoria}, e a memória de trabalho do controlador
        precisa sobrar para o programa de controle;
  \item o critério de fronteira é o \textbf{prazo}: o que precisa de resposta garantida fica no
        controlador; o resto sai.
\end{itemize}

\section{Qualidade do dado}

Indicador calculado sobre dado ruim é pior que nenhum indicador, porque produz confiança injustificada.
Quatro cuidados resolvem a maior parte dos problemas:

\begin{enumerate}
  \item \textbf{carimbo de tempo no ponto de origem}, não na chegada. Em rede não determinística, o
        instante da chegada não é o instante do evento;
  \item \textbf{qualidade acompanhando o valor}: distinguir ``zero'' de ``sem leitura''. São estados
        diferentes e exigem decisões diferentes;
  \item \textbf{resolução e período de amostragem coerentes} com o fenômeno
        (Capítulo~\ref{cap:es}). Amostrar devagar um evento rápido produz média enganosa;
  \item \textbf{unidade e escala documentadas} na lista de tags (Capítulo~\ref{cap:documentacao}).
        Duas áreas usando escalas diferentes para a mesma grandeza geram dois indicadores
        contraditórios --- e a discussão que ninguém quer ter.
\end{enumerate}

\begin{pratica}
Escolha uma máquina da sua planta ou bancada e construa um indicador de disponibilidade:
(a) liste as variáveis necessárias e diga se cada uma já existe ou precisa ser criada;
(b) defina as categorias de parada e quem vai classificá-las;
(c) monte o cálculo na Tabela~\ref{tab:indicadores};
(d) decida, pela Tabela~\ref{tab:onde-processar}, onde cada cálculo roda;
(e) escreva qual decisão será tomada com o resultado.
\end{pratica}

\begin{exercicios}
\begin{enumerate}[label=\alph*)]
  \item Dê um exemplo de dado, um de informação e um de indicador para uma linha de envase.
  \item Por que o OEE detalhado é mais útil para a engenharia que o OEE agregado?
  \item Uma máquina produziu 4.500 peças em 480 minutos programados, dos quais 40 foram de parada por
        falha e 20 de troca de produto planejada. Calcule a disponibilidade e diga que decisão a
        diferença entre as duas paradas permite tomar.
  \item Explique por que medir tempo de ciclo no supervisório pode esconder o problema que se quer
        investigar.
  \item Aplique a Tabela~\ref{tab:onde-processar} e decida onde rodar: (i) intertravamento de porta;
        (ii) média de consumo por turno; (iii) comparação de consumo entre meses; (iv) contagem de
        peças para o indicador de produção.
  \item Escreva a regra de fronteira entre CLP e borda em uma frase, e aplique-a a um caso seu.
  \item Descreva dois indicadores de degradação que se pode extrair apenas de contagem de ciclos e de
        tempo de etapa, sem sensor novo.
  \item Por que o carimbo de tempo deve ser gerado na origem? Dê um exemplo em que a diferença importa.
  \item Qual é a diferença prática entre ``valor zero'' e ``sem leitura'' em um indicador de consumo?
  \item Um colega propõe calcular o OEE de dez máquinas dentro de um dos CLPs. Aponte dois problemas
        técnicos com base no Capítulo~\ref{cap:memoria}.
\end{enumerate}
\end{exercicios}
